\documentclass[12pt]{article}
\usepackage[a4paper,margin=1in]{geometry}
\usepackage[T1]{fontenc}
\usepackage{lmodern,microtype}
\usepackage{amsmath,amssymb,amsthm,mathtools}
\usepackage{booktabs,array,enumitem,xcolor}
\usepackage[numbers,sort&compress]{natbib}
\usepackage[colorlinks=true,linkcolor=blue!55!black,citecolor=blue!55!black,urlcolor=blue!55!black]{hyperref}
\linespread{1.07}

\newtheorem{theorem}{Theorem}[section]
\newtheorem{proposition}[theorem]{Proposition}
\theoremstyle{remark}
\newtheorem{remark}[theorem]{Remark}
\newcommand{\norm}[1]{\left\lVert#1\right\rVert}

\title{Adjacent-Scale Transport of the Physical Edge Quartet\\
\large A Finite Obstruction to the Naive Haar Shell Map}
\author{Liang Wang}
\date{July 2026}

\begin{document}
\maketitle

\begin{abstract}
The outer four source-observable eigenmodes found in RH-200 form a canonical
finite packet at each of three noise scales.  This paper performs the first
direct adjacent-level transport test.  The packets at
$\sigma=0.04,0.02,0.01$ are put in common coordinates by the dyadic Haar
embedding.  We compare right and left invariant spaces, oblique Riesz
projectors, source and observation channel states, residues, eigenvalues,
and quartic characteristic coefficients.

The result is negative for the literal shell map.  Across four adjacent
scale/channel cases the largest right and left principal sines reach
$0.82175$ and $0.82388$; the largest relative oblique-projector defect is
$2.29068$.  The quartet-restricted intertwining defect reaches $0.65573$,
and one relative residue displacement reaches $8.18374$.  Even the finer
$0.02\to0.01$ step retains order-one projector defects.

The spectral correspondence itself does not disappear: every packet remains
conjugation closed and admits a unique minimum-cost matching.  Thus the
finite quartet is a valid local edge diagnostic but is not transported by
the unmodified coarse-cell embedding.  A renormalized map or a scalar
divisor formulation is required.  All results are finite floating audits;
no all-level shell, Fredholm determinant, Hilbert--P\'olya operator, or RH
conclusion is asserted.
\end{abstract}

\section{Problem and inherited packet}

RH-192--RH-201 corrected the state type of the temporal construction and
isolated a genuine source--observation packet \citep{WangRH201}.  At each
audited scale and on both physical channels, the four largest-modulus modes
are two nonreal conjugate pairs, all with nonzero transfer residue.  The
minimum radial gap after the fourth mode is $0.05949$ \citep{WangRH200}.

Local selection is not yet Gate A.  A spectral block can contribute to an
all-level determinant only if its identity is coherent under refinement.
The first candidate map is the isometric Haar embedding
\begin{equation}\label{eq:haar}
 J_ne_j=2^{-1/2}(e_{2j}+e_{2j+1}),
 \qquad J_n^*J_n=I.
\end{equation}
When matrix channel states double in both dimensions we also use a column
embedding $K_m$ and transport $X$ as $J_nXK_m^*$.

\section{Objects transported}

Let $A_c$ and $A_f$ denote adjacent coarse and fine physical operators.
Their outer quartets have right and left frames $V_c,W_c$ and $V_f,W_f$,
normalized by
\begin{equation}\label{eq:biorthogonal}
 W_c^*V_c=W_f^*V_f=I_4.
\end{equation}
The oblique packet projectors are
\begin{equation}\label{eq:projectors}
 P_c=V_cW_c^*,\qquad P_f=V_fW_f^*.
\end{equation}

The right-space transport angle is determined by the singular values of
$Q_f^*JQ_c$, where $Q_c,Q_f$ are orthonormal bases.  If $s_4$ is the smallest
singular value, we record
\begin{equation}\label{eq:angle}
 \sin\theta_{\max}=\sqrt{1-s_4^2}.
\end{equation}
Left spaces are treated independently because the operators are nonnormal.

For every simple mode $\lambda$, with projector $P_\lambda$, source $S$,
and observation $O$, the physical channel states are
\begin{equation}\label{eq:states}
 X_\lambda=P_\lambda S,
 \qquad Y_\lambda=P_\lambda^*O^*,
 \qquad r_\lambda=\langle Y_\lambda,X_\lambda\rangle_F.
\end{equation}
The comparison therefore tests both the base eigenspace and the actual
source--observation state.

\section{Predeclared correspondence}

At each endpoint we select the four largest-modulus eigenvalues before any
cross-level matching.  The minimum-total-distance permutation then matches
the two finite sets.  This permutation is used simultaneously for
eigenvalues, right/left modes, channel states, and residues.  It is not used
to alter the Haar map.

For the monic quartet polynomial
\begin{equation}\label{eq:quartic}
 D_\sigma(z)=\prod_{j=1}^4(z-\lambda_{\sigma,j}),
\end{equation}
we compare the complete coefficient vector.  This scalar comparison is
basis invariant and will become important after the state-space obstruction.

\section{Four adjacent-level cases}

The principal-angle and projector data are:
\begin{center}
\small
\begin{tabular}{llrrrr}
\toprule
step & side & right sine & left sine & $\norm{P_f-JP_cJ^*}_F/\norm{P_f}_F$ & polynomial error\\
\midrule
$0.04\to0.02$ & left  & $0.5766$ & $0.7654$ & $2.2907$ & $0.2402$\\
$0.04\to0.02$ & right & $0.7669$ & $0.6116$ & $2.2361$ & $0.2420$\\
$0.02\to0.01$ & left  & $0.8217$ & $0.6206$ & $1.0154$ & $0.3167$\\
$0.02\to0.01$ & right & $0.6267$ & $0.8239$ & $1.0077$ & $0.3096$\\
\bottomrule
\end{tabular}
\end{center}
None of these values is a small perturbative transport parameter.  In
particular, refining the grid does not make the naive projector defect small
on the two available transitions.

The operator and channel defects tell the same story.  The quartet-restricted
intertwining defects range from $0.23623$ to $0.65573$.  Source transport
defects lie between $0.78134$ and $0.81991$.  Observation defects are near
$2^{-1/2}$, reflecting both refinement and the changing physical
observation map.

\section{Mode and residue movement}

The largest eigenvalue displacement is $0.38160$ on the first transition and
$0.13643$ on the second.  Hence a recognizable spectral branch survives,
but its endpoint is moving.

Residues are less stable.  On the left $0.04\to0.02$ transition, one matched
mode changes from approximately
\begin{equation*}
 -0.09768-0.11776i
 \quad\text{to}\quad
 -0.00593-0.01569i.
\end{equation*}
Its relative displacement, normalized by the fine residue, is $8.18374$.
The minimum cosine between transported and fine source states over all
sixteen mode records is $0.48096$; the corresponding observation minimum is
$0.60264$.

These are gauge-invariant channel quantities.  Rephasing an eigenvector
cannot remove the residue discrepancy.

\section{Finite obstruction statement}

\begin{proposition}[Naive Haar shell obstruction]\label{prop:obstruction}
For the six physical endpoint packets and four adjacent cases audited here,
the rule
\begin{equation*}
 (P_c,X_c,Y_c)\longmapsto
 (JP_cJ^*,JX_cK^*,JY_cK^*)
\end{equation*}
does not provide a small-defect transport of the outer quartet.  In
particular, the maximum subspace sine exceeds $0.82$ and every relative
oblique-projector defect exceeds $1$.
\end{proposition}

The proposition is a statement about these matrices and this predeclared
map.  It is not an all-level theorem and does not rule out a scale-dependent
renormalization, a different cluster, or a scalar spectral construction.

\section{What survives the negative result}

Three facts remain useful:
\begin{enumerate}[leftmargin=2em]
\item the edge quartet is isolated and source observable at every endpoint;
\item the matched eigenvalues retain a common conjugate-branch geometry;
\item the monic quartic is meaningful without choosing eigenvector gauges.
\end{enumerate}
Thus the failure is specifically one of raw state transport.  It does not
invalidate the local Riesz packet or its spectral divisor.

\section{Audit protocol and dimensional ledger}

The left operator dimensions at $\sigma=0.04,0.02,0.01$ are respectively
$128,256,512$; the right dimensions are $64,128,256$.  Source column counts
are $64,128,256$ on both sides.  Thus every adjacent comparison doubles both
the base-state resolution and the source-column resolution.  The row and
column Haar maps are constructed independently with the normalization
\eqref{eq:haar}.

For reproducibility, the calculation uses the following fixed order:
\begin{enumerate}[leftmargin=2em]
\item construct each physical model directly from its declared $\sigma$;
\item compute the complete finite eigendecomposition;
\item select four indices by modulus before looking at another level;
\item biorthogonalize the selected right/left frames as one block;
\item solve one minimum-cost eigenvalue assignment;
\item reuse that assignment for every state, residue, and polynomial metric.
\end{enumerate}
No Procrustes correction is applied in RH-202.  This is essential: allowing
an endpoint-fitted unitary would answer a different question, developed only
in RH-205.

The order-one defects are far above eigensolver residual and conjugation
errors, but this does not make them interval statements.  Their role is to
reject the naive map as the next perturbative target, not to certify a sharp
lower bound for an exact continuum operator.

\section{Next mathematical wall}

The resolvent identity should be used to separate the causes of transport
failure.  If $E=A_fJ-JA_c$, then on a common resolvent set one expects
\begin{equation*}
 R_f(z)J-JR_c(z)=R_f(z)ER_c(z).
\end{equation*}
RH-203 develops this identity, integrates it over Riesz contours, and
separates projector defects from source defects.  Later layers must decide
whether to renormalize the state map or move first to the scalar divisor.

\section{Claim boundary}

The eigendecompositions and all norms in this paper are double-precision
finite calculations.  No interval enclosure of the input matrices,
resolvents, contours, or projectors is claimed.  The negative finite result
is robust as a diagnostic because the measured defects are order one, but
it is not promoted to an analytic zero-noise theorem.  Gate A remains open;
Gates B--E are untouched.

\bibliographystyle{plainnat}
\bibliography{references}
\end{document}
